\subsection{A targeted follow-up at 236 states}

The incoming 236 report already completes a plain-fold and one-neutral-helper
search on the new table. We add the freely synthesized two-helper
compensation method for the eight common-row pairs touching newly used
source or fold controls:
\[
 74/210,\ 185/229,\ 198/205,\ 198/219,\ 198/221,\ 198/246,
 \ 200/210,\ 210/232.
\]
The finite source profile has 28,926 radius-five windows. Old and new trace
horizons are 65, direct halt targets are allowed, and all halted heads/tapes
are identified. As before, the first helper must already make the merged
contexts pass before a second helper compensates its other uses. Across
7,552 retained-row trials, 54,553,565 first-helper trials and 4,715
second-helper trials, no 235-state candidate is found. This limited result
is not a state lower bound.

The independent audit also checks a positive composition fact: after the
four-row source rotation, the old $235\to98$ fold and new $223\to182$ fold
commute literally. Applying them in either order yields exactly the incoming
236 table. On the unrotated 237 table, the proposed $182/1$ three-step block
ends at state 185 instead of the required 212, isolating why the source
rotation is necessary. The small independent checker, exact inputs and
optional search recipes are under \texttt{fold236\_audit/}.

The new source rotation also commutes as a source rewrite with the previously
considered six-leaf NIP specialization, since their source prefixes are
disjoint. Their joint source still has 224 shared nodes and depth 49.
However, the attempted five-row NIP decoder graft still fails at the
unchanged $\mathrm{CI}_8$ instruction: the legal zero-register frame is
returned correctly in 636 steps by the new 236 table, whereas the altered
table returns in 648 steps with the wrong tape. Source rewrite compatibility
therefore does not justify composing the proposed binary patches.

